\section*{第 \ref{ch:b1:det} 章——行列式与线性方程组}

\begin{solution}{exo:b1:det:1}
$3 \times 2 - 1 \times 5 = 1$。

第二个：$L_2 \leftarrow L_2 - L_1$，$L_3 \leftarrow L_3 - L_2$（在
原始行上）得到行 $(1,2,3), (3,3,3), (3,3,3)$：两行相等，
行列式为 $0$。（萨吕法则验证：$45 + 84 + 96 - 105 - 48 -
72 = 0$。）

第三个：这是 $x = 1, 2, 3$ 时的\omterm{ex:b1:det:vandermonde}{范德蒙德行列式}
（见\cref{ex:b1:det:vandermonde}）：$(2-1)(3-1)(3-2) = 2$。
\end{solution}

\begin{solution}{exo:b1:det:2}
行列式等于（将所有列加到第一列，提取公因子）
$(\lambda + 2)$ 乘以
\[
\begin{vmatrix}
1 & 1 & \lambda\\ 1 & \lambda & 1\\ 1 & 1 & 1
\end{vmatrix}
= -(\lambda - 1)^2
\]
（用 $L_1 \leftarrow L_1 - L_3$，$L_2 \leftarrow L_2 - L_3$ 消元
并展开），得 $\det = -(\lambda+2)(\lambda-1)^2$。构成基
当且仅当 $\iff \det \neq 0 \iff \lambda \notin \{1, -2\}$。
\end{solution}

\begin{solution}{exo:b1:det:3}
第一个方程组：$\det = -7$；
$x = \frac{1}{-7}\begin{vmatrix} 5 & 1\\ 4 & -2\end{vmatrix}
= \frac{-14}{-7} = 2$,
$y = \frac{1}{-7}\begin{vmatrix} 2 & 5\\ 3 & 4\end{vmatrix}
= \frac{-7}{-7} = 1$。验算：$2(2) + 1 = 5$；$3(2) - 2 = 4$。

第二个方程组：作 $L_2 - L_1$ 与 $L_3 - 2L_1$ 后，各行变为
$(1,1,1)$、$(0,-2,0)$、$(0,-1,-3)$，于是
\[
\det A = \begin{vmatrix} 1&1&1\\ 1&-1&1\\ 2&1&-1\end{vmatrix}
= 1 \times \begin{vmatrix} -2 & 0\\ -1 & -3\end{vmatrix} = 6 .
\]
由克拉默法则，以 $(6,2,1)^{\mathsf T}$ 替换相应的列：
\[
x = \frac{6}{6} = 1, \qquad
y = \frac{12}{6} = 2, \qquad
z = \frac{18}{6} = 3
\]
（分子按同样方法算出）。验算：$1 + 2 + 3 = 6$；$1 -
2 + 3 = 2$；$2 + 2 - 3 = 1$。
\end{solution}

\begin{solution}{exo:b1:det:4}
对增广矩阵作行化简：$L_2 \leftarrow L_2 - 2L_1$，$L_3
\leftarrow L_3 - L_1$：
\[
\begin{pmatrix}
1 & 2 & -1 & 1 & 1\\
0 & 0 & 3 & -3 & 3\\
0 & 0 & 3 & -3 & 3
\end{pmatrix}
\to
\begin{pmatrix}
1 & 2 & -1 & 1 & 1\\
0 & 0 & 1 & -1 & 1\\
0 & 0 & 0 & 0 & 0
\end{pmatrix}.
\]。
主元未知量为 $x, z$；自由未知量为 $y, t$。回代得 $z =
1 + t$，$x = 1 - 2y + z - t = 2 - 2y$。解集为：
\[
\{(2 - 2y,\; y,\; 1 + t,\; t) : y, t \in \R\}
= (2, 0, 1, 0) + \operatorname{Vect}\bigl((-2,1,0,0),\,
(0,0,1,1)\bigr),
\]
这是一个仿射平面（维数 $2 = 4 - \operatorname{rk} 2$），位于 $\R^4$ 中。
\end{solution}

\begin{solution}{exo:b1:det:5}
用归纳法；$n = 1$ 时为空积 $= 1$。归纳步中，依次作
$C_k \leftarrow C_k - x_1 C_{k-1}$（$k = n, n-1, \dots, 2$，依此顺序，
使每次运算所用的列尚未被修改）。第一行
变为 $(1, 0, \dots, 0)$；在第 $i \geq 2$ 行中，第 $k$ 个元素
变为 $x_i^{k-1} - x_1 x_i^{k-2} = x_i^{k-2}(x_i - x_1)$。
按第一行展开，并从各行中提出因子 $(x_i - x_1)$，
对每行 $i$，得：
\[
V(x_1, \dots, x_n)
= \prod_{i=2}^{n} (x_i - x_1)\cdot V(x_2, \dots, x_n),
\]
再由归纳假设，得乘积 $\prod_{i<j}(x_j
- x_i)$。
\end{solution}

\begin{solution}{exo:b1:det:6}
按第一行展开 $D_n$：$D_n = 2 D_{n-1} -
1\cdot\begin{vmatrix} 1 & \ast\\ 0 & D_{n-2}\text{-块}
\end{vmatrix}$；第二个行列式按其第一
列展开，得 $D_{n-2}$。于是 $D_n = 2D_{n-1} - D_{n-2}$，即 $D_n -
D_{n-1} = D_{n-1} - D_{n-2}$：这些差为常数，其值为
$D_2 - D_1 = 1$。于是 $D_n = D_1 + (n - 1) = n + 1$。
（验算：$D_2 = 3$，$3\times3$ 的情形见
\cref{ex:b1:det:cofactor}：$D_3 = 4$。）
\end{solution}

\begin{solution}{exo:b1:det:7}
$C_1 \leftarrow C_1 + C_2 + C_3$ 使第一列成为常数列
$(m+2)$；提出该因子：
\[
\det = (m+2)\begin{vmatrix}
1 & 1 & m\\ 1 & m & 1\\ 1 & 1 & 1
\end{vmatrix}
\overset{L_1 - L_3,\ L_2 - L_3}{=}
(m+2)\begin{vmatrix}
0 & 0 & m-1\\ 0 & m-1 & 0\\ 1 & 1 & 1
\end{vmatrix}
= (m+2)\cdot\bigl(-(m-1)^2\bigr)
\]
（按第一列展开：元素 $1$ 所带的符号为
$+$，其余 $2 \times 2$ 行列式为 $0 \cdot 0 -
(m-1)(m-1) = -(m-1)^2$）。

讨论。$m \notin \{1, -2\}$：克拉默方程组；由方程组的对称性，
$x = y = z$，且每个方程都给出 $(m + 2)x = 1$：
唯一解 $\bigl(\frac{1}{m+2}, \frac{1}{m+2},
\frac{1}{m+2}\bigr)$。$m = 1$：三个方程都是 $x + y
+ z = 1$：解集为仿射平面 $x + y + z = 1$。$m = -2$：
将三个方程相加得 $0 = 3$：解集为空集。
\end{solution}

\begin{solution}{exo:b1:det:8}
($\Rightarrow$) 若 $A^{-1}$ 的元素均为整数：$\det A \cdot \det
A^{-1} = 1$，且两个行列式都是整数（由元素乘积之和得到）：
乘积为 $1$ 的两个整数都是 $\pm1$。

($\Leftarrow$) 伴随矩阵公式 $A^{-1} = \frac{1}{\det
A}\operatorname{Com}(A)^{\mathsf T}$（$n \leq 3$ 时可直接展开验证，
一般情形在本阶段不加证明地承认）中，
$\operatorname{Com}(A)$ 的元素均为整数（每个代数余子式都是整数行列式）；
除以 $\det A = \pm 1$ 仍得整数。
\end{solution}

\begin{solution}{exo:b1:det:9}
将所有列加到第一列：新第一列的每个元素为
$a + nb$；提出公因子，于是第一列全为一。然后
作行变换 $L_i \leftarrow L_i - L_1$（$i \geq 2$），
将左上角 $1$ 下方的元素全部消为零，这些行的对角线元素为 $a$，
其余位置为 $0$：矩阵为上三角矩阵，
对角线元素为 $(1, a, \dots, a)$。因此
\[
\det(aI + bJ) = (a + nb)\, a^{\,n-1} ,
\]
当且仅当 $a \neq 0$ 且 $a + nb \neq 0$ 时非零，即
\cref{exo:b1:matrices:9}中的条件。
\end{solution}

\begin{solution}{exo:b1:det:10}
分块运算（每一步都是 $n$ 次数量运算的复合，
由\cref{thm:b1:det:props}~(1) 允许）：
\[
\begin{vmatrix} A & B\\ B & A\end{vmatrix}
\overset{C_1 \leftarrow C_1 + C_2}{=}
\begin{vmatrix} A + B & B\\ A + B & A\end{vmatrix}
\overset{L_2 \leftarrow L_2 - L_1}{=}
\begin{vmatrix} A + B & B\\ 0 & A - B\end{vmatrix}
= \det(A+B)\,\det(A-B),
\]
最后一步用到分块三角矩阵的法则。
\end{solution}

\begin{solution}{exo:b1:det:11}
$C_1 \leftarrow C_1 + C_2 + C_3$ 使第一列成为常数列
$(a + b + c)$；提出公因子后，再作 $L_2 \leftarrow L_2 - L_1$，
$L_3 \leftarrow L_3 - L_1$：
\[
\Delta = (a+b+c)\begin{vmatrix}
1 & b & c\\
0 & a - b & b - c\\
0 & c - b & a - c
\end{vmatrix}
= (a+b+c)\bigl[(a-b)(a-c) + (b-c)^2\bigr],
\]
展开得 $(a-b)(a-c) + (b-c)^2 = a^2 + b^2 + c^2 - ab - bc
- ca$。在 $\C$ 中，取 $j^3 = 1$ 且 $1 + j + j^2 = 0$：
\begin{align*}
(a + jb + j^2c)(a + j^2b + jc)
&= a^2 + b^2 + c^2 + (j + j^2)(ab + bc + ca)\\
&= a^2 + b^2 + c^2 - ab - bc - ca ,
\end{align*}
由此得完整的因式分解。（从结构上看：列向量
$(1, j, j^2)^{\mathsf T}$ 满足 $M\,(1, j, j^2)^{\mathsf T}
= (a + jb + j^2c)(1, j, j^2)^{\mathsf T}$，同理对
$j^2$ 和 $1$：三个因子即循环矩阵的三个“特征值”，
这一思想将在本科二年级分册中系统化。）
\end{solution}

\begin{solution}{exo:b1:det:12}
记 $r = \operatorname{rk} A$。

\emph{存在可逆的 $r \times r$ 子矩阵。}任取 $r$ 个
$A$ 的\omterm{def:b1:vspaces:free}{线性无关}的列，记它们构成的矩阵为 $B \in \mathcal{M}_{n,r}$，
则 $\operatorname{rk} B = r$。由于行秩
等于列秩（\cref{thm:b1:matrices:rank}），$B$ 有 $r$
个\omterm{def:b1:vspaces:free}{线性无关}的行；保留这些行即得一个 $r \times r$ 子矩阵，
它是 $A$ 的子矩阵，其秩为 $r$，即可逆。

\emph{不存在更大的可逆子矩阵。}设 $S$ 为一个可逆的 $s \times s$
子矩阵，取自列 $j_1, \dots, j_s$ 与行 $i_1,
\dots, i_s$（均属于 $A$）。若 $\sum_k \lambda_k
C_{j_k} = 0$ 成立（所涉列均为相应的\emph{完整}列），
则只读取行 $i_1, \dots, i_s$，得 $\sum_k
\lambda_k S_k = 0$ 在 $S$ 的各列上成立，故所有 $\lambda_k = 0$
（$S$ 可逆）：列 $C_{j_1}, \dots, C_{j_s}$
为 $A$ 的列，\omterm{def:b1:vspaces:free}{线性无关}，且 $s \leq \operatorname{rk} A = r$。

因此 $\operatorname{rk} A$ 恰好是
可逆子矩阵的最大阶数。
\end{solution}

\begin{solution}{pb:b1:det:1}
\textbf{1.} $V(1,2,3,4) = (2-1)(3-1)(4-1)(3-2)(4-2)(4-3) = 1
\cdot 2\cdot 3\cdot 1\cdot 2\cdot 1 = 12$。插值问题
在互异节点上归结为求 $P$ 的系数，使其满足 $W c =
y$，其中 $W = (x_i^{\,j-1})$，且 $\det W = V \neq 0$：这是一个
克拉默方程组。

\textbf{2.} 从左至右逐列处理。$C_1$
是常数列 $P_0(x_i) = 1$（$P_0$ 为
$0$ 次\omterm{def:b1:poly:def}{首一多项式}）。假设第 $1, \dots, j-1$ 列已化为
纯幂列 $1, x_i, \dots, x_i^{\,j-2}$。由于 $P_{j-1}
= X^{j-1} + \sum_{k < j-1}\alpha_k X^k$，从 $C_j$ 中减去
组合 $\sum_k \alpha_k\,(\text{列 } x_i^k)$
——这一操作不改变行列式——得到
纯幂列 $x_i^{\,j-1}$。处理完最后一列后，
矩阵即为范德蒙德矩阵：$\det = V(x_1, \dots, x_n)$。

\textbf{3.} \omterm{def:b1:poly:def}{多项式} $(j-1)!\,B_{j-1}$ 均为\omterm{def:b1:poly:def}{首一}的，
次数为 $j - 1$，由第 2 题可得
\[
\det\bigl(B_{j-1}(m_i)\bigr)
= \frac{V(m_1, \dots, m_n)}{0!\,1!\cdots(n-1)!} .
\]
左端为一个元素均为\emph{整数}的矩阵的行列式，
（$B_k$ 在 $\Z$ 上取整数值：见
\cref{pb:b1:vspaces:1}中的第 16--17 题），故为整数；
且其为正，因为 $V(m_1, \dots, m_n) > 0$，只要 $m_1 <
\dots < m_n$。于是超阶乘\omterm{def:b1:arith:divides}{整除}
全体两两差之积。

\textbf{4.} 提出各行的公因子 $x_i$（第 $i$ 行）：
$\det(x_i^{\,j})_{j = 1..n} = x_1\cdots x_n\,
\det(x_i^{\,j-1}) = x_1\cdots x_n\,V$.

\textbf{5.} $(W^{\mathsf T}W)_{ij} = \sum_k x_k^{\,i-1}
x_k^{\,j-1} = p_{i+j-2}$，记 $S = W^{\mathsf T}W$。于是 $\det S =
\det(W^{\mathsf T})\det W = V^2$
（由\cref{thm:b1:det:props}~(2),(4)）。对实数 $x_i$，$\det S = V^2
\geq 0$，且 $S$ 可逆当且仅当 $V \neq 0$，即
$x_i$ 两两互异——一种仅
由幂和即可算出的符号判别法。

\textbf{6.} 插值条件 $\sum_{k}
c_k\,x_i^{\,k} = y_i$ 构成方程组 $Wc = y$，由 $\det W = V \neq
0$ 立即得到存在性与唯一性。这是本书中的第三种
证明：显式公式见
\cref{thm:b1:poly:lagrange}，核的论证见
\cref{ex:b1:linmaps:interpolation}，此处用克拉默法则。

\textbf{7.} 由克拉默法则，$c_{n-1} = \det W'/\det W$，其中 $W'$ 为把 $W$
的最后一列换为 $y$ 所得的矩阵。沿
该列展开 $\det W'$：
\[
\det W' = \sum_{i=1}^n (-1)^{i+n} y_i\,V(x_1, \dots, \widehat{x_i},
\dots, x_n) .
\]
注意到 $V = V(\setminus i)\cdot\prod_{j<i}(x_i - x_j)\prod_{j>i}
(x_j - x_i)$，而把第二个乘积改写会带来
因子 $(-1)^{n-i}$：
\[
(-1)^{i+n}\,\frac{V(\setminus i)}{V}
= \frac{(-1)^{i+n}(-1)^{n-i}}{\prod_{j\neq i}(x_i - x_j)}
= \frac{1}{\prod_{j\neq i}(x_i - x_j)} ,
\]
于是 $c_{n-1} = \sum_i y_i/\prod_{j \neq i}(x_i - x_j)$——
即又是\omterm{pb:b1:vspaces:1}{差商}公式。

\textbf{8.} 作 $L_3 \leftarrow L_3 - L_1$，得各行为 $(1, x_1,
x_1^2)$，$(0, 1, 2x_1)$，$(0,\ x_2 - x_1,\ (x_2-x_1)(x_2+x_1))$；
按第一列展开，并提出公因子 $(x_2 - x_1)$：
\[
(x_2 - x_1)\begin{vmatrix} 1 & 2x_1\\ 1 & x_2 + x_1
\end{vmatrix} = (x_2 - x_1)(x_2 - x_1) = (x_2 - x_1)^2 .
\]
当 $x_1 \neq x_2$ 时该行列式非零：以下
$P(x_1) = u$，$P'(x_1) = v$，$P(x_2) = w$ 这组条件关于
$P \in \R_2[X]$ 的系数的\omterm{def:b1:det:system}{线性方程组}可用克拉默法则求解——含
二重节点的 埃尔米特插值是适定的。

\textbf{9.} 设 $P = a + bX + cX^2$，其中 $a = P(0) = 1$，$b = P'(0)
= 0$，$a + b + c = P(1) = 2$：$c = 1$，故 $P = 1 + X^2$，唯一。
相容性检验：此处 $x_1 = 0$，$x_2 = 1$，第 8 题中的行列式
为 $(1 - 0)^2 = 1 \neq 0$。

\textbf{10.} 直接计算：
\[
C_2 = \frac{1}{(a_1+b_1)(a_2+b_2)} - \frac{1}{(a_1+b_2)(a_2+b_1)}
= \frac{(a_1+b_2)(a_2+b_1) - (a_1+b_1)(a_2+b_2)}
{\prod_{i,j}(a_i+b_j)} ,
\]
而分子展开为 $a_1b_1 + a_2b_2 - a_1b_2 - a_2b_1
= (a_2 - a_1)(b_2 - b_1)$：即差比和。

\textbf{11.} 对 $i < n$，行 $i$ 的新元素为
\[
\frac{1}{a_i + b_j} - \frac{1}{a_n + b_j}
= \frac{a_n - a_i}{(a_i + b_j)(a_n + b_j)} .
\]
提出 $(a_n - a_i)$（对各行 $i < n$），再
提出 $\frac1{a_n + b_j}$（对各列 $j$）：余下的
矩阵的元素为 $\frac1{a_i + b_j}$（行 $i < n$），以及常数 $1$（行 $n$）
——即矩阵 $M$，带上所述的前置因子。

\textbf{12.} 对 $M$，当 $j < n$ 时作运算 $C_j \leftarrow
C_j - C_n$，使第 $n$ 行变为 $(0, \dots, 0, 1)$，且在行 $i
< n$，
\[
\frac{1}{a_i + b_j} - \frac{1}{a_i + b_n}
= \frac{b_n - b_j}{(a_i + b_j)(a_i + b_n)} .
\]
提出 $(b_n - b_j)$（对各列 $j < n$），并提出 $\frac1{a_i +
b_n}$（对各行 $i < n$），再沿最后一行展开
（符号 $(-1)^{n+n} = +1$）：剩下的行列式为
$C_{n-1}$。综合第 11--12 题中的各因子：
\[
C_n = \frac{\prod_{i<n}(a_n - a_i)\,\prod_{j<n}(b_n - b_j)}
{\prod_{j}(a_n + b_j)\,\prod_{i<n}(a_i + b_n)}\;C_{n-1},
\]
由归纳法（以 $C_1 = \frac1{a_1+b_1}$ 为基础）
恰好拼出\omterm{pb:b1:det:1}{柯西行列式}：因子 $(a_j - a_i)(b_j
- b_i)$ 遍取所有数对，分母则遍取所有和 $(a_i + b_j)$。

\textbf{13.} 该公式为零当且仅当存在 $a_j = a_i$ 或 $b_j =
b_i$：柯西矩阵可逆当且仅当两组数各自
两两互异。\omterm{pb:b1:det:1}{希尔伯特矩阵}取 $a_i = i$，$b_j = j - 1$。对 $n =
2$：分子 $(2-1)(1-0) = 1$，分母 $1\cdot2\cdot2\cdot3
= 12$，$\det H_2 = \frac1{12}$。对 $n = 3$，分子为
$\bigl[(1)(2)(1)\bigr]^2 = 4$，分母为 $(1\cdot2\cdot3)
(2\cdot3\cdot4)(3\cdot4\cdot5) = 6\cdot24\cdot60 = 8640$：
$\det H_3 = \frac{4}{8640} = \frac1{2160}$。当 $n = 2$ 时其逆为：
\[
H_2^{-1} = 12\begin{pmatrix} \frac13 & -\frac12\\[2pt]
-\frac12 & 1\end{pmatrix}
= \begin{pmatrix} 4 & -6\\ -6 & 12 \end{pmatrix},
\]
全为整数（这一现象对任意 $n$ 都成立）。

\textbf{14.} 方程组的系数矩阵是柯西矩阵，
由第 13 题知，当 $b_j$（及 $a_i$）
两两不同时它可逆，解唯一。解释：具有单\omterm{def:b1:fractions:field}{极点}的有理
函数 $R = \sum_j \frac{c_j}{X + b_j}$ 由其
$n$ 个值 $R(a_1), \dots, R(a_n)$ 唯一确定，反过来
任一此类数据组恰好对应一个这样的有理函数——
这正是部分分式的存在性与
唯一性（\cref{thm:b1:fractions:complex}）的抽样形式。

\textbf{15.} 对 $X^3 + pX + q$ 用韦达定理：$\lambda_1 + \lambda_2 +
\lambda_3 = 0$，$\sum_{i<j}\lambda_i\lambda_j = p$，于是 $p_1 = 0$
从而 $p_2 = p_1^2 - 2p = -2p$。每个根满足 $\lambda^3 =
-p\lambda - q$。求和得 $p_3 = -p\,p_1 - 3q = -3q$。两边乘以
$\lambda$ 再求和：$p_4 = -p\,p_2 - q\,p_1 = 2p^2$。

\textbf{16.} 由第 5 题（恒等式 $S = W^{\mathsf T}W$
及 $\det S = V^2$ 在 $\C$ 上成立），
\[
V^2 = \begin{vmatrix}
3 & 0 & -2p\\
0 & -2p & -3q\\
-2p & -3q & 2p^2
\end{vmatrix}
= 3\bigl(-4p^3 - 9q^2\bigr) + (-2p)\bigl(0 - 4p^2\bigr)
= -4p^3 - 27q^2 ,
\]
按第一行展开。

\textbf{17.} 重根即两个 $\lambda_i$ 相等，即
$V = 0$，即 $\operatorname{disc} = -4p^3 - 27q^2 = 0$。对
$X^3 - 3X + 2$：$4(-3)^3 + 27\cdot4 = -108 + 108 = 0$，
与二重根 $1$ 相符，该二重根来自 $(X-1)^2(X+2)$。

\textbf{18.} 实三次\omterm{def:b1:poly:def}{多项式}的非实根以共轭对的形式出现，
故当 $\operatorname{disc} \neq
0$ 时恰有两种情形。三个互异实根：$V$ 为非零实数，故
$\operatorname{disc} = V^2 > 0$。一个实根 $\lambda_1$ 与
$\lambda_3 = \conj{\lambda_2} \notin \R$，则
\[
(\lambda_2 - \lambda_1)(\lambda_3 - \lambda_1) =
\abs{\lambda_2 - \lambda_1}^2 > 0,
\qquad
\lambda_3 - \lambda_2 = -2\iu\,\operatorname{Im}\lambda_2 \neq 0,
\]
于是 $V$ 为非零纯虚数，且
$\operatorname{disc} = V^2 < 0$。两种符号刻画了
两种情形。

\textbf{19.} 在双交错式中取 $b_i = a_i$，
则分子为 $\prod_{i<j}(a_j - a_i)^2 > 0$，分母
为 $\prod_{i,j}(a_i + a_j) > 0$（元素皆为正）：该
行列式为正。（用后面的语言说：核
$\frac1{x+y}$ 是正定的。）

\textbf{20.} 由第 2--3 题，该行列式等于
$V(2,4,7)/(0!\,1!\,2!) = \frac{(4-2)(7-2)(7-4)}{2} =
\frac{30}{2} = 15$。直接可得矩阵为
\[
\begin{pmatrix}
1 & 2 & 1\\
1 & 4 & 6\\
1 & 7 & 21
\end{pmatrix},
\qquad
\det = (84 - 42) - 2(21 - 6) + (7 - 4) = 42 - 30 + 3 = 15 .
\]

\textbf{21.} 设 $\sum_i c_i\,(\lambda_i^{\,k})_{k} = 0$ 作为
数列。依次取 $k = 0, 1, \dots, n-1$，得 $W^{\mathsf
T}c = 0$，其中 $W = (\lambda_i^{\,j-1})$ 可逆（$\det = V
\neq 0$，$\lambda_i$ 互异）：$c = 0$。几何数列
\omterm{def:b1:vspaces:free}{线性无关}。

\textbf{22.} 取 $a = b = (1, 2, 3)$，分子为 $\bigl[(2-1)(3-1)
(3-2)\bigr]^2 = 4$；分母 $\prod_{i,j}(i + j) =
(2\cdot3\cdot4)(3\cdot4\cdot5)(4\cdot5\cdot6) = 24\cdot60\cdot120
= 172800$。于是 $\det\bigl(\frac1{i+j}\bigr) = \frac{4}{172800}
= \frac1{43200}$。

\textbf{23.} 若 $x_i = x_j$，交换两变量保持该点不变，
但 $F$ 须变号：$F = -F$，故 $F = 0$
在该点处。\omterm{def:b1:arith:divides}{整除性}：将 $F$ 视为单个
变量 $x_n$ 的\omterm{def:b1:poly:def}{多项式}，系数属于其余变量；它
在 $n - 1$ 个“值” $x_1, \dots, x_{n-1}$ 处为零，于是
由\cref{thm:b1:poly:factor}逐次因式分解，得 $F =
\prod_{i<n}(x_n - x_i)\cdot G$，其中 $G$ 为\omterm{def:b1:poly:def}{多项式}。
前置因子在交换两指标 $i, j < n$ 时不变，故
$G$ 关于 $x_1, \dots, x_{n-1}$ 交错，归纳法
即完成：$\prod_{i<j}(x_j - x_i)$ \omterm{def:b1:arith:divides}{整除} $F$。

\textbf{24.} $D = \det(x_i^{\,j-1})$ 是关于
$x_i$ 的\omterm{def:b1:poly:def}{多项式}；交换两变量即交换两行，故 $D$ 是
交错的，由第 23 题，$D = c\,\prod_{i<j}(x_j - x_i)$
对某个\omterm{def:b1:poly:def}{多项式} $c$ 成立。总次数：$D$ 的次数 $\leq 0 +
1 + \dots + (n-1) = \binom n2$，而乘积的次数恰为
$\binom n2$：$c$ 为常数。单项式 $x_2\,x_3^2\cdots
x_n^{\,n-1}$ 的系数为 $1$（$D$ 中对角线乘积），而
在乘积中亦为 $1$（每个
因子取指标较大的变量）：$c = 1$，范德蒙德公式
无需归纳即可得出。

\textbf{25.} (i) 交错性——公理“两列相等
则行列式为零”——是核心：它产生了每个
因子 $(x_j - x_i)$、$(a_j - a_i)$、$(b_j - b_i)$，均在
题目中出现。（ii）恒等式 $\det S = V^2$ 替换了
单独看为复数、无法直接求出的根，换为它们的幂和，
幂和是系数的实\omterm{def:b1:poly:def}{多项式}，因此互异性
化为可计算实数的符号。（iii）\omterm{ex:b1:det:vandermonde}{范德蒙德
行列式}支配\omterm{def:b1:poly:def}{多项式}插值，
\omterm{pb:b1:det:1}{柯西行列式}支配部分分式及抽样的
有理函数（\omterm{pb:b1:det:1}{希尔伯特矩阵}是其中最著名的
特例）。（iv）任一交错\omterm{def:b1:poly:def}{多项式}都可被
$\prod_{i<j}(x_j - x_i)$ \omterm{def:b1:arith:divides}{整除}，再由次数计数可把此类
\omterm{def:b1:poly:def}{多项式}确定至多相差一个常数——这正是该乘积
反复出现于行列式因
重合而为零之处。第三部分的定理是\emph{柯西
双交错式}。
\end{solution}
